\begin{description}
\item[(T1.1)] \textbf{True.} \hfill(2 points)

  The colour of the clear sky during night is the same as during daytime,
  since the spectrum of sunlight reflected by the Moon is almost the same as
  the spectrum of sunlight. Only the intensity is lower.

\item[(T1.2)] \textbf{True.} \hfill(2 points)

  If the Earth's axis were perpendicular to its orbital plane, the celestial
  equator will coincide with ecliptic and the Sun will remain along the
  celestial equator every day. However, as the Earth's orbit is elliptical,
  the true sun would still lead or lag mean sun by a few minutes on different
  days of year.

\item[(T1.3)] \textbf{False.} \hfill(2 points)

  The semi-major axis of the orbit of the body will be less than that of
  Uranus. However the minor body's orbit may have a high eccentricity, in
  which case it may go outside that of Uranus.

\item[(T1.4)] \textbf{False.} \hfill(2 points)

  The centre of mass of Sun--Jupiter pair is just outside the Sun. Thus, if
  all gas giants are on same side of the Sun, the centre of mass of Solar
  system is definitely outside the Sun.

\item[(T1.5)] \textbf{True.} \hfill(2 points)

  For photons the wavelength increases when the Universe expands.
\end{description}
